\chapter{Os Quatro Subespaços Fundamentais}
\label{cap:subespacos-fundamentais}

\textbf{Importante:} este capítulo é baseado em \cite{slide4doyuri} e em \cite{axler2026}.

Neste capítulo, nós iremos desenvolver mais profundamente os conceitos relacionados a espaços vetoriais, e apresentar os quatro subespaços fundamentais de uma matriz. Além disso, nós iremos enunciar um teorema muito importante, que os relaciona entre si: o Teorema do Posto (que será generalizado no próximo capítulo para toda transformação linear).

\section{Posto de uma Matriz}

No capítulo anterior, nós definimos um número $r$ como o número de pivôs de uma matriz $A$. Agora, nós formalizaremos esta noção, e a utilizaremos para falar sobre o número de soluções de um sistema linear.

\begin{Definition}[Posto]
	Seja $A$ uma matriz $m\times n$ qualquer com $r$ pivôs. Ao número $r$ se dá o nome de \textbf{posto} (ou \textbf{rank}) da matriz $A$.
\end{Definition}

Uma consequência imediata desta definição é a seguinte desigualdade:

\begin{Proposition}
	Seja $A$ uma matriz $m\times n$ de posto $r$. Então:
	
	$$
	r\leq\min\{m, n\}.
	$$
\end{Proposition}

Para os casos em que ocorre uma igualdade, a matriz $A$ recebe nomes especiais com relação ao posto:

\begin{Definition}[Posto completo nas colunas]
	Seja $A$ uma matriz $m\times n$ de posto $r$. Se $r = n$, nós dizemos que $A$ tem \textbf{posto completo nas colunas}.
\end{Definition}

\begin{Definition}[Posto completo nas linhas]
	Seja $A$ uma matriz $m\times n$ de posto $r$. Se $r = m$, nós dizemos que $A$ tem \textbf{posto completo nas linhas}.
\end{Definition}

É interessante analisar as consequências de uma matriz $A$ possuir posto completo nas colunas ou nas linhas. Se $A$ possui posto completo nas colunas, então todas as colunas são pivôs, e o sistema possui $0$ ou $1$ solução, dependendo de $b$. Por sua vez, se $A$ possui posto completo nas linhas, então todas as linhas são pivôs, e o sistema possui infinitas soluções, independente de $b$.

Também é interessante analisar a relação entre o posto e a forma escalonada reduzida (com permutações) de $A$. De fato, dado que o processo de escalonamento e, depois, as permutações em $R$, revelam uma certa estrutura, é esperado que tal estrutura varie conforme a completude do posto.

Se $A$ é uma matriz quadrada com posto completo (ou seja, $r = n = m$), então:

$$
R' = I
$$.
Consequentemente, neste caso, $N(A) = \{0\}$ e $C(A) = \mathbb{R}^n$.

Se $A$ é uma matriz $m\times n$ com posto completo nas colunas (ou seja, $r = n < m$), então:

$$
R' = \begin{bmatrix}
	I_{n\times n} \\
	0_{(m - n)\times n}
\end{bmatrix}
$$
Consequentemente, neste caso, $N(A) = \{0\}$ e as linhas de zeros restringem $C(A)$.

Se $A$ é uma matriz $m\times n$ com posto completo nas linhas (ou seja, $r = m < n$), então:

$$
R' = \begin{bmatrix}
	I_{m\times m} & F_{m\times(n - m)}
\end{bmatrix}
$$
Consequentemente, neste caso, $N' = \begin{bmatrix}
	-F \\
	I
\end{bmatrix}$ e $C(A) = \mathbb{R}^m$.

No caso geral, se $A$ é uma matriz $m\times n$, com $r < m$ e $r < n$, então:

$$
R' = \begin{bmatrix}
	I_{r\times r} & F_{r\times(n - r)} \\
	0_{(m - r)\times r} & 0_{(m - r)\times(n - r)}
\end{bmatrix}
$$
Consequentemente, neste caso, $N' = \begin{bmatrix}
	-F \\
	I
\end{bmatrix}$ e $C(A)$ é restringido pelas linhas de zeros de $R'$.

\section{Os Quatro Subespaços Fundamentais}

Tendo definido o conceito de posto, nós estamos prontos para dizer quais são os quatro subespaços fundamentais de uma matriz. Dois deles nós já conhecemos: são o espaço coluna e o núcleo. Os outros dois nós definiremos a seguir:

\begin{Definition}[Quatro subespaços fundamentais]
	Seja $A$ uma matriz $m\times n$ qualquer. Então, os \textbf{quatro subespaços fundamentais} associados a $A$ são:
	
	$\mathbf{1)}$ $C(A) = \{Ax\in\mathbb{R}^m; x\in\mathbb{R}^n\}$ (espaço coluna);
	
	$\mathbf{2)}$ $N(A) = \{x\in\mathbb{R}^n; Ax = 0\}$ (núcleo)
	
	$\mathbf{3)}$ $C(A^T) = \{A^Ty\in\mathbb{R}^n; y\in\mathbb{R}^m\}$ (espaço linha);
	
	$\mathbf{4)}$ $N(A^T) = \{y\in\mathbb{R}^m; A^Ty = 0\}$ (núcleo à esquerda).
\end{Definition}

Eles são chamados de subespaços fundamentais porque, juntos, descrevem toda a estrutura de uma matriz, e também porque estão relacionados entre si de uma maneira que nós ainda veremos.

Uma pequena observação sobre o núcleo à esquerda é o fato de que $A^Ty = (y^TA)^T$. Sabemos que $A^Ty$ é uma combinação linear das linhas de $A$. Assim sendo, $A^Ty = 0 \iff y^TA = 0$. Por isso, $N(A^T)$ é chamado de núcleo à esquerda.

Muito em breve, nós desenvolveremos muitos conceitos e teoremas que se baseiam fortemente nos quatro subespaços fundamentais de uma matriz. No entanto, para que nós possamos chegar a este ponto, precisamos dar um passo atrás e entender alguns conceitos que são comuns a todos os espaços e subespaços vetoriais. Nós faremos isto nas próximas seções, e depois retomaremos a discussão sobre os subespaços fundamentais.

\section{Base de um Espaço Vetorial}

Quando nós discutimos o conceito de posto, nós estávamos preocupados em entender qual é a relação do número de pivôs de uma matriz com o número de soluções de um sistema linear associado a ela. Porém, o que exatamente produz um pivô? Ou melhor, o que exatamente produz a falta de um pivô?

Vamos analisar um exemplo simples. Considere a matriz $A = \begin{bmatrix}
	1 & 3 \\
	2 & 6
\end{bmatrix}$. Ao fazer o processo de eliminação, nós obtemos $U = \begin{bmatrix}
	1 & 3 \\
	0 & 0
\end{bmatrix}$, que possui apenas um pivô. Nós podemos perceber que a segunda coluna é uma combinação linear da primeira (e que o mesmo acontece nas linhas). Agora, altere a matriz $A$, obtendo $A = \begin{bmatrix}
	1 & 3 \\
	2 & 5
\end{bmatrix}$. Ao fazer o processo de eliminação, nós obtemos $U = \begin{bmatrix}
	1 & 3 \\
	0 & -1
\end{bmatrix}$, que possui dois pivôs. Nós podemos, por sua vez, que a segunda coluna de $A$ não é uma combinação linear da primeira (e o mesmo vale para as linhas).

Assim sendo, algo nos sugere que a relação entre um vetor ser ou não uma combinação linear de outros é algo importante. É precisamente isto que nós desenvolveremos agora.

\begin{Definition}[Independência linear]
	Seja $V$ um espaço vetorial. Diremos que $v_1, v_2, \cdots, v_n\in V$ são \textbf{linearmente independentes (LI)} se
	
	$$
	\alpha_1v_1 + \alpha_2v_2 + \cdots + \alpha_nv_n = 0 \iff \alpha_1 = \alpha_2 = \cdots = \alpha_n = 0.
	$$
\end{Definition}

No caso da segunda matriz do nosso exemplo, é possível demonstrar que as suas colunas são linearmente independentes. Vamos fazê-lo:

\begin{Proposition}
	Sejam $v_1 = \begin{bmatrix}
		1 \\
		2
	\end{bmatrix}$ e $v_2 = \begin{bmatrix}
		3 \\
		5
	\end{bmatrix}$ em $\mathbb{R}^2$. Então, $v_1$ e $v_2$ são linearmente independentes.
\end{Proposition}

\begin{Proof}
	Considere a equação $\alpha_1v_1 + \alpha_2v_2 = 0$. Esta equação é, na verdade, um sistema de equações, dado por:
	
	$$
	\begin{cases}
		\alpha_1 + 3\alpha_2 = 0 \\
		2\alpha_1 + 5\alpha_2 = 0
	\end{cases}
	$$
	Realizando a eliminação, nós obtemos:
	
	$$
	\begin{cases}
		\alpha_1 + 3\alpha_2 = 0 \\
		-\alpha_2 = 0
	\end{cases}
	$$
	Este sistema possui uma única solução: $\alpha_1 = 0$ e $\alpha_2 = 0$. Logo, $v_1$ e $v_2$ são linearmente independentes.
	
	\qed
\end{Proof}

Por sua vez, se determinados vetores em um espaço não são linearmente independentes, eles são linearmente dependentes.

\begin{Definition}[Dependência linear]
	Seja $V$ um espaço vetorial. Diremos que $v_1, v_2, \cdots, v_n$ são \textbf{linearmente dependentes (LD)} se existem $\alpha_1, \alpha_2, \cdots, \alpha_n\in\mathbb{R}$, nem todos nulos, tais que:
	
	$$
	\alpha_1v_1 + \alpha_2v_2 + \cdots + \alpha_nv_n = 0.
	$$
\end{Definition}

Um exemplo de vetores linearmente dependentes é dado pelas colunas da primeira matriz do nosso exemplo. De fato, $\begin{bmatrix}
	3 \\
	6
\end{bmatrix} = 3\begin{bmatrix}
 	1 \\
 	2
\end{bmatrix}$.

No capítulo anterior, nós definimos o conjunto \textit{span} como sendo o conjunto de todas as combinações lineares de um certo grupo de vetores. Assim sendo, nós podemos definir o que significa um conjunto de vetores gerar um espaço em termos do conjunto span:

\begin{Definition}[Gerador]
	Sejam $v_1, v_2, \cdots, v_n$ vetores tais que $\operatorname{span}(\{v_1, v_2, \cdots, v_n\}) = V$. Então, dizemos que $v_1, v_2, \cdots, v_n$ \textbf{geram} $V$.
\end{Definition}

Na prática, isto quer dizer que qualquer vetor $w\in V$ pode ser escrito como uma combinação linear de $v_1, v_2, \cdot, v_n$. No entanto, note que nós não dissemos que $v_1, v)_2, \cdots, v_n$ são LI. De fato, isto não é uma condição necessária para que um grupo de vetores gere um espaço ou um subespaço vetorial. Os geradores podem ser LD. Porém, muitas vezes nós estamos interessados no caso em que os vetores geradores são LI. Eles são o menor grupo de vetores que gera um espaço, e isso é algo especial. Assim sendo, nós podemos definir o conceito de base:

\begin{Definition}[Base de um espaço vetorial]
	O conjunto de vetores $\{v_1, v_2, \cdots, v_n\}$ é dito uma \textbf{base} para $V$ se:
	
	$\mathbf{1)}$ $v_1, v_2, \cdots, v_n$ são linearmente independentes;
	
	$\mathbf{2)}$ $v_1, v_2, \cdots, v_n$ geram $V$.
\end{Definition}

